\chapter{Predicting a new structure}
\label{ch:inference}

\begin{lstlisting}
poraque-inference new_structure/ \
    --output predictions/new_structure
\end{lstlisting}

The directory needs only \file{POSCAR}, \file{INCAR} and \file{POTCAR}. The
script computes the external potential, predicts the density, then the kinetic
energy density, and writes all three in \file{CHGCAR} format.

\section{Choosing the grid}

Resolved in this order, first hit wins:

\begin{enumerate}
  \item \opt{--grid NX NY NZ} --- explicit;
  \item \opt{--like FILE} --- adopt an existing volumetric file's grid, which is
    what you want when comparing against a reference calculation;
  \item \opt{--resolution N} --- longest axis, preserving the aspect ratio;
  \item otherwise \opt{--encut} (default \textbf{200\,eV}) through the usual
    \opt{ENCUT}/\opt{PREC} rule.
\end{enumerate}

The default cutoff is Poraqu\^e's own, not the one the reference calculation
happened to use: a genuinely new structure has no prior calculation to inherit
from, and the same geometry must give the same grid either way. When the two
differ the log says so.

\begin{ptip}
200\,eV is chosen to land near the grid the shipped models were trained on. On
the reference cells it gives $42^3$ (27-atom) and $40^3$ (32-atom), against a
training \opt{resolution} of $32$ --- close enough that the operator is
interpolating rather than extrapolating. Raising it to 400\,eV would give
$60^3$, a substantially finer mesh than anything the models have seen.
\end{ptip}

\begin{pwarn}
A Fourier neural operator is resolution-\emph{flexible}, not
resolution-\emph{indifferent}. Evaluating far from the grid density a model was
trained on is extrapolation, and no metric printed here can detect it. The
script warns when the requested grid departs markedly from the training
resolution; pass \opt{--resolution 32} to evaluate exactly where the models
were fitted.
\end{pwarn}

\section{Comparing against a reference}

\begin{lstlisting}
poraque-inference run/ \
    --like run/CHGCAR --compare --plot-dir results/plots/check
\end{lstlisting}

\opt{--compare} scores the prediction against any reference files present,
bringing them onto the prediction grid by spectral truncation. \opt{--plot-dir}
adds cross-section and parity figures.

\section{Sanity checks the script performs}

\begin{itemize}
  \item the predicted electron count against $\sum_a Z^{\mathrm{val}}_a$;
  \item the number of negative density voxels, if any;
  \item the bound $\tau \ge \tau_{\rm vW}[\rho]$ on the chained output, with the
    minimum margin.
\end{itemize}

A large electron-count error or a violated bound means the prediction should
not be trusted, whatever the field looks like.

\section{Visualising}

All outputs are \file{CHGCAR}-format and open directly in VESTA. Use
\opt{--write-chg} for an additional coarse-formatted \file{CHG}.
